
En
effet, on a d\'ej\`a not\'e l'inclusion @@M0@@. Pour l'inclusion en sens inverse, on note que @@M1@@ est dans @@M2@@ si et seulement si la suite des @@M3@@ (@@M4@@) est quasi-compacte\,; comme @@M5@@ est quasi-compact, il suffit pour ceci que la suite des @@M6@@ ait la m\^eme propri\'et\'e\,; d'o\`u l'assertion.

\textbf{Lemme 4.3.} \emph{Soient @@M7@@, @@M8@@ propres sur un sch\'ema @@M9@@ et @@M10@@ un morphisme surjectif (resp. le morphisme d'inclusion du sch\'ema des z\'eros d'une section @@M11@@ d'un faisceau ample @@M12@@, sous les hypothèses de 3.8). Alors}
@@M13@@

En effet, l'assertion est cons\'equence formelle de 4.2, (XII 1.2), 3.1, et 3.5 (resp. 3.8).

\textbf{4.4.} Soient @@M14@@ un corps alg\'ebriquement clos, @@M15@@ un @@M16@@-sch\'ema propre, et @@M17@@ un faisceau inversible sur @@M18@@. On dit que @@M19@@ est \uline{alg\'ebriquement \'equivalent} \`a @@M20@@ (resp. @@M21@@-\uline{\'equivalent} \`a @@M22@@) s'il existe un @@M23@@-sch\'ema connexe de type fini @@M24@@, un faisceau inversible @@M25@@ sur @@M26@@, et deux @@M27@@-points @@M28@@ de @@M29@@ tels que @@M30@@ et @@M31@@ (resp. s'il existe un entier non nul @@M32@@ tel que @@M33@@ soit alg\'ebriquement \'equivalent \`a @@M34@@), ou ce qui revient facilement au m\^eme, si le @@M35@@-point @@M36@@ de @@M37@@ qui correspond \`a @@M38@@, est dans @@M39@@ (resp. dans @@M40@@).

Rappelons aussi que @@M41@@ est dit \uline{num\'eriquement \'equivalent} \`a @@M42@@ si pour toute courbe int\`egre ferm\'ee @@M43@@ dans @@M44@@, le nombre d'intersection @@M45@@ est @@M46@@ (voir 2.1 ou X 4.5).

